\section{Closed subspaces of affine schemes}
\label{section-closed-in-affine}

\begin{example}
\label{example-closed-immersion-affines}
Let $R$ be a ring.
Let $I \subset R$ be an ideal.
Consider the morphism of affine schemes
$i : Z = \Spec(R/I) \to \Spec(R) = X$.
By Algebra, Lemma \ref{algebra-lemma-spec-closed}
this is a homeomorphism of $Z$ onto a closed
subset of $X$. Moreover, if $I \subset \mathfrak p \subset R$
is a prime corresponding to a point $x = i(z)$, $x \in X$,
$z \in Z$, then on stalks we get the map
$$
\mathcal{O}_{X, x} = R_{\mathfrak p}
\longrightarrow
R_{\mathfrak p}/IR_{\mathfrak p} = \mathcal{O}_{Z, z}
$$
Thus we see that $i$ is a closed immersion of locally ringed spaces,
see Definition \ref{definition-closed-immersion-locally-ringed-spaces}.
Clearly, this is (isomorphic) to the closed subspace associated
to the quasi-coherent sheaf of ideals $\widetilde I$, as in
Example \ref{example-closed-subspace}.
\end{example}

\begin{lemma}
\label{lemma-closed-immersion-affine-case}
\begin{slogan}
For affine schemes, closed immersions correspond to ideals.
\end{slogan}
Let $(X, \mathcal{O}_X) = (\Spec(R), \mathcal{O}_{\Spec(R)})$
be an affine scheme. Let $i : Z \to X$ be any closed immersion
of locally ringed spaces. Then there exists a unique ideal
$I \subset R$ such that the morphism $i : Z \to X$ can be identified
with the closed immersion $\Spec(R/I) \to \Spec(R)$
constructed in Example \ref{example-closed-immersion-affines} above.
\end{lemma}

\begin{proof}
This is kind of silly! Namely, by Lemma \ref{lemma-closed-immersion}
we can identify $Z \to X$ with the closed subspace associated to
a sheaf of ideals $\mathcal{I} \subset \mathcal{O}_X$ as in
Definition \ref{definition-closed-subspace} and
Example \ref{example-closed-subspace}.
By our conventions this sheaf of ideals is locally generated
by sections as a sheaf of $\mathcal{O}_X$-modules.
Hence the quotient sheaf $\mathcal{O}_X / \mathcal{I}$
is locally on $X$ the cokernel
of a map $\bigoplus_{j \in J} \mathcal{O}_U \to \mathcal{O}_U$.
Thus by definition, $\mathcal{O}_X / \mathcal{I}$ is quasi-coherent.
By our results in Section \ref{section-quasi-coherent-affine}
it is of the form $\widetilde S$ for some $R$-module $S$.
Moreover, since $\mathcal{O}_X = \widetilde R \to \widetilde S$
is surjective we see by Lemma \ref{lemma-extension-quasi-coherent}
that also $\mathcal{I}$ is quasi-coherent, say $\mathcal{I} = \widetilde I$.
Of course $I \subset R$ and $S = R/I$ and everything is clear.
\end{proof}













